% !TEX root = ../Main.tex
\chapter{角平分线}

角平分线题的关键是\textbf{添辅助线}，而辅助线的核心思想是凑出\textbf{相似（或全等）}。下面先看两种添线方式，再用面积比证角平分线定理。

\section{几何辅助线构造}

设 $OP$ 平分 $\angle O$，$P$ 在角平分线上。

\textbf{课内做法（最常见）}：从 $P$ 向角的两边各作垂线，垂足为 $M,N$。由角平分线上的点到两边距离相等，$PM=PN$，于是 $\triangle OPM\cong\triangle OPN$。

\begin{center}
\begin{tikzpicture}[>=Stealth,scale=1.0]
  \coordinate (O) at (0,0);
  \draw[thick] (O) -- (4,1.4);
  \draw[thick] (O) -- (4,-1.4);
  \draw[thick] (O) -- (3.6,0);
  \coordinate (P) at (2.2,0);
  \fill (P) circle (1.1pt) node[below]{$P$};
  \fill (O) circle (1.1pt) node[left]{$O$};
  % perpendicular feet
  \coordinate (M) at ({2.2*cos(19.29)*cos(19.29)},{2.2*cos(19.29)*sin(19.29)});
  \coordinate (N) at ({2.2*cos(19.29)*cos(19.29)},{-2.2*cos(19.29)*sin(19.29)});
  \draw[dashed] (P) -- (M) node[above right]{\footnotesize$M$};
  \draw[dashed] (P) -- (N) node[below right]{\footnotesize$N$};
\end{tikzpicture}
\end{center}

\textbf{课外做法}：过 $P$ 作\textbf{垂直于角平分线 $OP$} 的直线，交两边于 $S,T$。因为 $OP$ 既平分 $\angle O$、又垂直于 $ST$，所以 $\triangle OST$ 是等腰三角形，$P$ 是 $ST$ 的中点（$SP=PT$）。

\begin{center}
\begin{tikzpicture}[>=Stealth,scale=1.0]
  \coordinate (O) at (0,0);
  \draw[thick] (O) -- (4,1.4);
  \draw[thick] (O) -- (4,-1.4);
  \draw[thick] (O) -- (3.6,0);
  \coordinate (P) at (2.2,0);
  \coordinate (S) at (2.2,0.77);
  \coordinate (T) at (2.2,-0.77);
  \draw[dashed] (S) -- (T);
  \fill (P) circle (1.1pt) node[below right]{$P$};
  \fill (O) circle (1.1pt) node[left]{$O$};
  \fill (S) circle (1.1pt) node[above]{$S$};
  \fill (T) circle (1.1pt) node[below]{$T$};
\end{tikzpicture}
\end{center}

\rmk{两种做法都把问题拉回到\textbf{相似（全等）}上——剖析到相似关系，许多角平分线问题就一通百通了。}

\section{角平分线定理}

\thm{角平分线定理}{
    在 $\triangle ABC$ 中，$AD$ 平分 $\angle BAC$，$D$ 在 $BC$ 上，则
    \[
    \frac{AB}{BD}=\frac{AC}{CD}.
    \]
}

\begin{center}
\begin{tikzpicture}[>=Stealth,scale=1.0]
  \coordinate (A) at (0.8,2.6);
  \coordinate (B) at (-1.4,0);
  \coordinate (C) at (4,0);
  \coordinate (D) at ($(B)!{2.97/(2.97+3.83)}!(C)$);
  \draw[thick] (A) -- (B) -- (C) -- cycle;
  \draw[thick,red!70!black] (A) -- (D);
  \node at (0.0,1.0) {$S_1$};
  \node at (1.2,1.0) {$S_2$};
  \fill (A) circle (1.3pt) node[above]{$A$};
  \fill (B) circle (1.3pt) node[below left]{$B$};
  \fill (C) circle (1.3pt) node[below right]{$C$};
  \fill (D) circle (1.3pt) node[below]{$D$};
\end{tikzpicture}
\end{center}

\pf[面积之比]{
    记 $S_1=S_{\triangle ABD}$，$S_2=S_{\triangle ACD}$，并设 $\angle BAD=\angle CAD=\alpha$。\textbf{同一个面积比，用两种方式数。}

    其一，两个三角形都以 $A$ 为公共顶点、夹角都是 $\alpha$：
    \[
    \frac{S_1}{S_2}=\frac{\frac12\,AB\cdot AD\sin\alpha}{\frac12\,AC\cdot AD\sin\alpha}=\frac{AB}{AC}.
    \]

    其二，$\triangle ABD$ 与 $\triangle ACD$ 以 $BD,DC$ 为底，公用从 $A$ 到 $BC$ 的同一条高 $h$：
    \[
    \frac{S_1}{S_2}=\frac{\frac12\,BD\cdot h}{\frac12\,DC\cdot h}=\frac{BD}{DC}.
    \]

    两式左端相同，故 $\dfrac{AB}{AC}=\dfrac{BD}{DC}$，即 $\dfrac{AB}{BD}=\dfrac{AC}{CD}$。
}
